\section{半导体进口供应链风险指数构建}

为综合评价不同产品的进口来源结构风险，本文构建半导体进口供应链风险指数 SIRI。该指数按产品--年份计算，分值越高表示该产品在该年份的进口来源结构风险越高。SIRI 包含来源集中风险、政策暴露风险、替代不足风险和结构波动风险四个维度。

设 $s_{i,p,t}$ 为年份 $t$ 中国自来源国 $i$ 进口产品 $p$ 的份额。来源集中风险采用 HHI 衡量：
\[
C_{p,t}=\sum_i s_{i,p,t}^{2}.
\]
政策暴露风险使用美国来源份额衡量，即 $P_{p,t}=s_{US,p,t}$。替代不足风险使用非美国来源内部 HHI 衡量，若非美国来源越集中，则说明替代空间越弱。结构波动风险使用相邻年份来源份额的总变差距离：
\[
V_{p,t}=\frac{1}{2}\sum_i |s_{i,p,t}-s_{i,p,t-1}|.
\]
四个维度先在全部产品--年份样本内进行 0--1 标准化，再等权加总：
\[
SIRI_{p,t}=100\times(0.25C_{p,t}^{*}+0.25P_{p,t}^{*}+0.25A_{p,t}^{*}+0.25V_{p,t}^{*}).
\]
其中星号表示标准化后的指标。等权设定可以降低主观赋权影响。为检验结论稳健性，本文进一步将政策暴露风险权重提高至 0.40，其余三项权重设为 0.20。

\begin{figure}[H]
    \centering
    \includegraphics[width=0.92\textwidth]{siri_trend_by_product_v02.png}
    \caption{半导体相关产品 SIRI 风险指数变化趋势}
    \label{fig:siri-trend}
\end{figure}

表 \ref{tab:siri-ranking} 显示 2024 年 SIRI 排序。854239 和 854232 的综合风险相对更高，主要与来源集中和结构暴露有关。848620 虽然美国份额下降明显，但其 2024 年综合风险排序低于两类集成电路产品，说明设备类和集成电路类风险来源并不相同。

\begin{table}[H]
\centering
\caption{2024年 SIRI 风险指数排序}
\label{tab:siri-ranking}
\begin{tabular}{llrrr}
\toprule
排序 & HS6编码 & 基准SIRI & 政策暴露加权SIRI & 排序变化 \\
\midrule
1 & 854239 & 39.47 & 33.84 & 0 \\
2 & 854232 & 34.99 & 27.99 & 0 \\
3 & 848620 & 18.16 & 20.32 & 0 \\
4 & 854231 & 14.15 & 17.03 & 0 \\
\bottomrule
\end{tabular}
\end{table}


表 \ref{tab:siri-sensitivity} 给出权重敏感性检验。将政策暴露风险权重提高后，四类产品排序未发生变化，说明本文关于相对风险排序的结论不完全依赖单一权重设定。

\begin{table}[H]
\centering
\caption{SIRI 权重敏感性检验}
\label{tab:siri-sensitivity}
\begin{tabular}{lrrrrr}
\toprule
HS6编码 & 基准排序 & 政策加权排序 & 排序变化 & 基准得分 & 政策加权得分 \\
\midrule
854239 & 1 & 1 & 0 & 39.47 & 33.84 \\
854232 & 2 & 2 & 0 & 34.99 & 27.99 \\
848620 & 3 & 3 & 0 & 18.16 & 20.32 \\
854231 & 4 & 4 & 0 & 14.15 & 17.03 \\
\bottomrule
\end{tabular}
\end{table}

